\par\medskip
\begin{center}\hypertarget{native-164-article118}{}\pdfbookmark[2]{118.}{article-118}118.\end{center}
Ist der einer Beobachtung entsprechende Abstand des Himmelskörpers von der Erde bereits annähernd bekannt, so lässt sich die Beobachtung auf mehrere Arten von der Wirkung der \emph{Aberration} befreien. Sie beruhen auf den verschiedenen in Artikel 71 angegebenen Verfahren. Es sei $t$ die wirkliche Beobachtungszeit und $\theta$ die Zeitspanne, in der das Licht vom Himmelskörper zur Erde gelangt. Man erhält sie, indem man $493^{\mathrm{s}}$ mit dem Abstand multipliziert. Ferner sei $l$ der beobachtete Ort und $l'$ derselbe Ort, der mithilfe der täglichen geozentrischen Bewegung auf die Zeit $t+\theta$ reduziert ist. Es sei \SourceLeaf{165}$l''$ der Ort $l$, befreit von dem Teil der Aberration, den die Planeten mit den Fixsternen gemeinsam haben. Es sei $L$ der wahre Ort der Erde zur Zeit $t$ (also der Tafelort, vergrößert um $20''{,}25$) und schließlich ${}'L$ der wahre Ort der Erde zur Zeit $t-\theta$. Dann gilt
\par I. $l$ ist der wahre Ort des Himmelskörpers, gesehen von ${}'L$ aus zur Zeit $t-\theta$
\par II. $l'$ ist der wahre Ort des Himmelskörpers, gesehen von $L$ aus zur Zeit $t$
\par III. $l''$ ist der wahre Ort des Himmelskörpers, gesehen von $L$ aus zur Zeit $t-\theta$.
\par\noindent Im Verfahren I behält man also den beobachteten Ort bei und setzt an die Stelle der wirklichen Zeit die fiktive Zeit $t-\theta$, für die man auch den Ort der Erde berechnet. Im Verfahren II ändert man nur die Beobachtung. Diese Änderung erfordert jedoch außer dem Abstand auch die tägliche Bewegung. Im Verfahren III bringt man an der Beobachtung eine vom Abstand unabhängige Korrektur an und ersetzt die wirkliche Zeit durch die fiktive Zeit $t-\theta$. Dabei behält man den Ort der Erde bei, der der wirklichen Zeit entspricht. Das erste Verfahren ist weitaus am bequemsten, wenn der Abstand bereits so weit bekannt ist, dass man die Zeitkorrektur $\theta$ mit hinreichender Genauigkeit berechnen kann.

Ist dieser Abstand jedoch noch völlig unbekannt, so lässt sich keines dieser Verfahren unmittelbar anwenden. Beim ersten Verfahren hat man zwar den geozentrischen Ort des Himmelskörpers, doch fehlen die Zeit und der Ort der Erde, die vom unbekannten Abstand abhängen. Beim zweiten hat man diese beiden Größen, aber der geozentrische Ort fehlt. Beim dritten hat man den geozentrischen Ort des Himmelskörpers und den Ort der Erde, doch fehlt die zu diesen Angaben gehörige Zeit.

Wie sollen wir also bei unserer Aufgabe vorgehen, wenn in einem solchen Fall eine auch hinsichtlich der Aberration genaue Lösung verlangt wird? Am einfachsten ist es, die Bahn zunächst unter Vernachlässigung der Aberration zu bestimmen. Da die Aberration niemals eine beträchtliche Wirkung hervorrufen kann, ergeben sich daraus die Abstände gewiss mit solcher Genauigkeit, dass man die Beobachtungen nach einem der eben beschriebenen Verfahren von der Aberration befreien und die Bahnbestimmung genauer wiederholen kann. Hierbei ist das dritte Verfahren den übrigen bei Weitem vorzuziehen. Beim ersten muss man alle vom Ort der Erde abhängigen Rechenoperationen von vorn beginnen. Beim zweiten (das nur anwendbar ist, wenn genügend Beobachtungen vorliegen, um daraus die tägliche Bewegung abzuleiten) müssen alle vom geozentrischen Ort des Himmelskörpers abhängigen Rechenoperationen erneut ausgeführt werden. Beim dritten lassen sich hingegen, sofern schon die erste Rechnung auf geozentrischen Orten beruhte, die von der Aberration der Fixsterne befreit waren, alle vorbereitenden Rechenoperationen, die vom \SourceLeaf{166}Ort der Erde und vom geozentrischen Ort des Himmelskörpers abhängen, unverändert in die neue Rechnung übernehmen. Auf diese Weise lässt sich die Aberration sogar schon in die erste Rechnung einbeziehen, wenn das zur Bahnbestimmung verwendete Verfahren so beschaffen ist, dass sich die Werte der Abstände ergeben, bevor die korrigierten Zeiten in die Rechnung eingeführt werden müssen. Wegen der Aberration allein ist dann keine zweite Rechnung nötig, wie bei der ausführlicheren Behandlung unserer Aufgabe deutlicher werden wird.
